\documentclass{article}
\usepackage{graphicx}
\usepackage[margin=1.5cm]{geometry}
\usepackage{amsmath}

\begin{document}
\twocolumn

\title{Friday Warm Up: Unit 0, Coulomb's Law and E-fields}
\author{Prof. Jordan C. Hanson}

\maketitle

\section{Memory Bank}

\begin{itemize}
\item $\vec{F}_{12} = k q_1 q_2 / r^2 ~~\hat{r}$ ... Coulomb force.
\item $k = 8.988 \times 10^{9}$ N C$^{-2}$ m$^{2}$ ... The constant of proportionality in Coulomb force.
\item $\vec{F} = q\vec{E}$ ... Definition of an electric field.
\item $1e = 1.602 \times 10^{-19}$ Coulombs ... Fundamental unit of charge.
\end{itemize}

\section{Charge, Coulomb Force}

\begin{enumerate}
\item Suppose you have a charge of +1 nC, or $10^{-9}$ Coulombs, separated from another identical charge.  (a) What will be the force between them? (b) Will the charges accelerate toward each other or away from each other? \\ \vspace{4cm}
\item How many protons are required to create a \textit{total charge} of +1 nC, or $10^{-9}$ Coulombs? \\ \vspace{3cm}
\item If two charges are a distance $r$ apart, and then $r$ is tripled, by what factor does the force decrease?
\begin{itemize}
\item A: 3
\item B: 6
\item C: 9
\item D: 12
\end{itemize}
\end{enumerate}

\section{Mathematics Review}

\begin{enumerate}
\item Compute the dot-product, $\vec{a} \cdot \vec{b}$.
\begin{itemize}
\item $\vec{a} = 10\hat{i} + 10\hat{j}$, $\vec{b} = 10\hat{i} + 10\hat{j}$.
\end{itemize}
\item Compute the cross-product, $\vec{a} \times \vec{b}$.
\begin{itemize}
\item $\vec{a} = -2\hat{i} + 2\hat{j}$, $\vec{b} = 4\hat{i} + 4\hat{j}$.
\item $\vec{a} = 2\hat{i} - 1\hat{j}$, $\vec{b} = -2\hat{i} - 3\hat{j}$.
\end{itemize}
\end{enumerate} \vspace{4cm}

\section{Electrostatic Equilibrium}

\begin{enumerate}
\item Two small balls, each of mass 5.0 g, are attached to silk threads 50 cm long, which are in turn tied to the same point on the ceiling. When the balls are given the same charge Q, the threads hang at 5.0 degrees to the vertical. What is the magnitude of Q? What are the signs of the two charges?
\end{enumerate}

\end{document}
